% Insert the first subsection in Section III and the second in Section IV.
% Draft text. Report held-out outcomes only after the experiment is run.
\subsection{Revision-scoped Dependency Repair}
Let \(G^0=(V,E_s^0,E_r,E_o)\) be the authored graph and
\(d_r=(u,k,v)\) the edit at revision \(r\), with anchor shot \(u\),
edited attribute key \(k\), and new value \(v\). A candidate
\(c=(r,u,w,k,h_w)\) binds a downstream shot \(w\) to the input digest
\(h_w\) from the last published revision. We consider only forward
state edges to shots with a compatible declared attribute binding.
For each candidate we record a verdict in
\(\{\mathrm{Confirmed},\mathrm{Rejected},\mathrm{Unknown}\}\)
together with evidence-frame digests. Confirmed candidates contribute
keyed state edges \((u,w,k)\) to the effective graph of revision \(r\).
Rejected candidates contribute no edge. Authored edges remain in force.
A conservative fallback may schedule work for an Unknown candidate
without recording its identity as Confirmed.

\subsection{Evidence Acquisition and Revision-Safe Execution}
For each sampled frame, the identity probe records host and part visibility,
marker shape, color, location, and the same-instance relation. The prompt
asks for observations about the host and part. The decision rule compares
these fields with a prespecified marker ontology. A visible part linked to
a matching host marker confirms the candidate. Rejection requires at least
one explicitly excluded identity, complete coverage of the designated
frames, and no unresolved frame or confirmed target. Missing visibility,
uncertain ownership, and malformed evidence remain Unknown. This rule
limits exclusion to the evidence supplied by the designated frames.

Let \(Q_r\) be the shots included by an explicit conservative fallback,
and let \(G_r\) denote the effective graph with confirmed repairs.
The regeneration set is the DAIR closure from
\(\{u\}\cup Q_r\) under the edited key on \(G_r\). With the current
runtime, each accepted repair propagates from the newly affected child
over the authored graph. This is sufficient for the supported
anchor-to-child repairs. Under a traversal of the applicable edges, one
closure evaluation costs \(O(|V|+|E_r^{\mathrm{eff}}|)\), where
\(E_r^{\mathrm{eff}}\) denotes all effective edges at revision \(r\).
Processing \(m\) candidates incrementally can require up to \(m\)
such traversals, in addition to at most \(Bm\) model calls at budget
\(B\) per candidate. The prototype uses two ordered pairs of frames
and stops early on positive linked evidence.

The runtime accepts a resolution only for the current revision and its
pinned input digest. Publication waits for all pinned candidates to be
resolved or given an explicit fallback, and for the required current
outputs to commit. This protects revision consistency. The shot-selection
guarantee is conditional on the correctness and coverage of the supplied
evidence.

For a future cost-sensitive policy, one may compare reuse, regeneration,
and acquisition of another frame using
\(J(a)=\lambda_1 C_{\mathrm{render}}(a)+\lambda_2R_{\mathrm{miss}}(a)
+\lambda_3R_{\mathrm{false}}(a)+\lambda_4C_{\mathrm{evidence}}(a)\).
Estimating those risks requires disjoint calibration data. The present
structured probe evaluates a fixed evidence budget and preserves Unknown
candidates for conservative fallback.
